% =====================================================================
%  Claimed, Not Achieved -- arXiv manuscript source
% =====================================================================
%  Build:
%      pdflatex main && bibtex main && pdflatex main && pdflatex main
%  or  latexmk -pdf main
%
%  NOTHING IN THIS FILE CARRIES A TRANSCRIBED FIGURE. Every number is
%  either a macro from tables/macros.tex or a row body in tables/*.tex,
%  all generated by `python3 eval/paper_tables.py` from the ledger. If a
%  number here looks stale, regenerate -- do not correct it by hand,
%  because a hand-edit is a number the ledger no longer backs.
%
%  Document class is deliberately plain `article`: the protocol fixes the
%  venue only after the evidence, so this stays venue-neutral and gets
%  reskinned into a template later. Both figures are drawn with base
%  LaTeX rules and tabulars rather than TikZ, so the source compiles on a
%  minimal TeX installation as well as on arXiv.
%
%  NOTE ON THE ARXIV SUBMISSION FORM: the abstract below is the PDF
%  abstract. arXiv's metadata abstract field is capped at 1920
%  characters; this one exceeds it, so the web form needs a trimmed
%  version. Trim from the calibration sentence onward -- do not trim the
%  claimed-versus-delivered sentence, which is the paper.
% =====================================================================
\documentclass[11pt]{article}

\usepackage[T1]{fontenc}
\usepackage[utf8]{inputenc}
\usepackage[letterpaper,margin=1.1in]{geometry}
\usepackage{amsmath}
\usepackage{amssymb}
\usepackage{array}
\usepackage{booktabs}
\usepackage{graphicx}
\usepackage{xcolor}
\usepackage{enumitem}
\usepackage{microtype}
\usepackage[numbers,sort&compress]{natbib}
\usepackage[hidelinks,breaklinks]{hyperref}   % load last

% ---- palette, shared with the HTML article -------------------------
\definecolor{verified}{HTML}{2C6A4A}
\definecolor{falsecol}{HTML}{9E4127}
\definecolor{neutral}{HTML}{78868D}

% ---- generated numbers ---------------------------------------------
%% GENERATED by `python3 eval/paper_tables.py` -- do not edit.
%% Regenerate rather than correct: a hand-edit here is a number the
%% ledger no longer backs.
\newcommand{\NPrimary}{356}
\newcommand{\NArms}{7}
\newcommand{\OnN}{52}
\newcommand{\OffN}{51}
\newcommand{\OnSucc}{67.3\%}
\newcommand{\OffSucc}{29.4\%}
\newcommand{\OnClaim}{71.2\%}
\newcommand{\OffClaim}{74.5\%}
\newcommand{\OnFC}{3.8\%}
\newcommand{\OffFC}{45.1\%}
\newcommand{\SuccDiff}{$+$37.9pp}
\newcommand{\SuccCI}{[$+$20.2, $+$56.9]}
\newcommand{\FCDiff}{$-$41.3pp}
\newcommand{\FCCI}{[$-$60.8, $-$23.1]}
\newcommand{\Resamples}{10,000}%


% Read a generated block of tabular rows.
%
% This must use the TeX primitive, not LaTeX's \input. LaTeX's \input
% expands to `\@@input <file> \relax`, and that trailing \relax arrives
% exactly where TeX is peeking for the start of the next row. \relax is
% neither \noalign nor \omit, so TeX opens a fresh row and inserts the
% column template -- and the \bottomrule that follows then sits inside a
% cell, failing with "! Misplaced \noalign". The primitive takes a
% space-terminated filename and leaves nothing behind.
\makeatletter
\newcommand{\inputrows}[1]{\@@input #1 }
\makeatother

% ---- a run-in head: the atomic unit of this manuscript -------------
\newcommand{\runin}[1]{\paragraph{#1}}

% ---- callout blocks, base LaTeX only -------------------------------
\newenvironment{callout}[1]
  {\par\medskip\noindent\rule{\linewidth}{0.9pt}\par\nobreak\smallskip
   \noindent{\footnotesize\textsc{#1}}\par\nobreak\smallskip\small}
  {\par\smallskip\noindent\rule{\linewidth}{0.4pt}\par\medskip}

% ---- horizontal bars for Figure 1 ----------------------------------
\newlength{\barmax}
\setlength{\barmax}{0.46\linewidth}
\newcommand{\pbar}[2]{\textcolor{#2}{\rule[-0.08ex]{#1\barmax}{0.85ex}}}

\title{\textbf{Claimed, Not Achieved}\\[.35em]
  \large A controlled same-model ablation of a deterministic harness,\\
  \large in an open world where establishing ground truth is the
  engineering problem}

% TODO before submission: authors, affiliations, corresponding address.
\author{\textcolor{falsecol}{[authors and affiliation to be set]}}
\date{Draft manuscript, \today{} --- circulated for comment, not submitted}

\begin{document}
\maketitle

\begin{abstract}
\noindent
Open-world LLM agents are increasingly evaluated, and increasingly
trained, on a completion signal the agent itself produces. This is a
defensible engineering choice --- in an unbounded action space there is
no answer key --- but it makes the party with an interest in the number
the party who certifies it. The resulting error has been measured at
scale, and every such measurement was taken in a domain where ground
truth is a lookup. In this paper we measure it where verification is
itself the engineering problem. We ablate a five-layer deterministic
harness under a fixed model, world, task set and commit, layer by layer,
recording three signals that are never collapsed: the agent's own claim,
an independent scorer's verdict, and the taxonomy derived from their
disagreement. Across \NPrimary{} pre-registered attempts --- \NArms{}
arms $\times$ 13 tasks $\times$ 4 replicates --- the harnessed and
unharnessed agents claim success at statistically indistinguishable
rates, \OnClaim{} against \OffClaim{}, while delivering \OnSucc{}
against \OffSucc{} verified success. Ablation raises false success from
\OnFC{} to \OffFC{} (\FCDiff{}, 95\% CI \FCCI{}): the agent overstates
on 23 of \OffN{} attempts and understates on none. We then locate the
mechanism and manipulate it --- false successes concentrate where the
starting kit already contained the goal item, and stripping only that
item drives them from 18/24 to 0/24 under a decision rule fixed in
advance. A pre-registered per-layer ablation finds \emph{no} single
layer responsible, and the stratified data explains why: the layers are
redundant on this failure mode, so each is individually undetectable by
construction. Because the paper's subject is that self-report cannot be
trusted, we hold our own instrument to the same standard --- a blind
human pass drawn \emph{by scorer verdict} and sealed under a published
hash before any evidence card was rendered certifies the false-success
verdict at 20/21, Wilson $[77.3, 99.2]$, and declines to certify the
structure-scoring path rather than averaging it away. Pre-registration,
deviations log, ledger, evidence files and analysis code are released.
\end{abstract}

% =====================================================================
\begin{figure}[t]
\centering
\begin{tabular}{@{}r@{\hspace{1em}}l@{\hspace{1em}}r@{}}
\multicolumn{3}{@{}l@{}}{\footnotesize\textsc{what the agent said}}\\[.3em]
\footnotesize full harness & \pbar{0.712}{neutral} & \footnotesize \OnClaim{}\\
\footnotesize no harness   & \pbar{0.745}{neutral} & \footnotesize \OffClaim{}\\[.9em]
\multicolumn{3}{@{}l@{}}{\footnotesize\textsc{what the world showed}}\\[.3em]
\footnotesize full harness & \pbar{0.673}{verified} & \footnotesize \OnSucc{}\\
\footnotesize no harness   & \pbar{0.294}{falsecol} & \footnotesize \OffSucc{}\\
\end{tabular}
\caption{\textbf{The whole paper in four bars.} Same model, same world,
same commit; the harness is the only difference. The two arms are
indistinguishable on the signal an evaluation pipeline usually records
and differ by \SuccDiff{} on the signal it is meant to stand for. The
harness does not make the agent claim less --- it makes the claims true.}
\label{fig:teaser}
\end{figure}

% =====================================================================
\section{Introduction}
\label{sec:intro}

Open-ended embodied agents built on large language models have moved
quickly from demonstrations to systems that accumulate their own
competence. Voyager explores a Minecraft world indefinitely and writes
reusable skills into a permanent library \citep{voyager}; Luban builds
creative artifacts through autonomous embodied verification
\citep{luban}; MineEvolve distils successful executions into accumulated
knowledge across long horizons \citep{mineevolve}; MineExplorer
evaluates open-world exploration with rule-based milestones
\citep{mineexplorer}. The common ambition is long-horizon autonomy in a
world nobody enumerated in advance, and the common ingredient is a loop
that decides, on its own, when a task is finished.

That the loop grades itself is not an oversight, and the case for it
deserves stating. In a domain with an unbounded action space there is no
answer key to consult: the set of world states that count as ``a
shelter'' cannot be written down, and instrumenting an environment to
recognise each one is a per-task engineering cost that scales with the
curriculum. The agent's own report, by contrast, is free, universal, and
available for every task --- including tasks invented at run time. A
critic agent generalises where a checker does not. Self-report became
the default in open worlds because it is the only completion signal that
is always there.

The cost is specific, and it compounds. In Voyager the critic's approval
is what admits code to the permanent skill library, so a single wrong
judgment contaminates the library and the error is re-executed on every
reuse \citep{voyager}. Luban's pipeline rests on the same self-issued
verification \citep{luban}, and MineEvolve distils executions its own
loop judged successful \citep{mineevolve}. In each case the
certification sits upstream of everything downstream --- what gets
stored, what gets retried, what gets published as a success rate --- and
it is issued by the party with an interest in the answer.

That this concern is real rather than hypothetical is already
established, and we claim neither part of it. That agents assert
completion they have not achieved has a name --- \emph{false success}
--- and has been measured at scale: across 9,876 tau2-bench and 1,879
AppWorld trajectories it accounts for 45--48\% of failures in
single-control domains and 75.8\% among self-assessing coding
trajectories \citep{falsesuccess}. That capability and self-report are
separable axes is also established: VIGIL separates world completion
from self-termination across 20 models and 1,000 frozen episodes and
finds systems with near-identical execution differing by up to 19.7
points in benchmark success \citep{vigil}. What is missing is elsewhere,
and it is a claim about \emph{measurement regimes}. Every measurement
above was taken where establishing ground truth is cheap. tau2-bench
checks a database row; AppWorld checks an API result; SpreadsheetBench
compares a cell to an exact match. Those settings establish that the
phenomenon exists and is common. They say nothing about the setting
where verification is \emph{itself the engineering problem} --- which is
the setting anything acting in a physical world operates in.

\begin{table}[t]
\centering\small
\begin{tabular}{@{}lll@{}}
\toprule
 & API / tool benchmarks & open world\\
\midrule
action space   & enumerable, documented    & effectively unbounded\\
ground truth   & look up a row, run a test & must independently read world state\\
horizon        & short, few dependencies   & long, tool and material chains\\
agent's priors & docs describe affordances & affordances must be discovered\\
failure        & terminal                  & normal; recovering is part of the task\\
\bottomrule
\end{tabular}
\caption{Two regimes, and what changes between them.}
\label{tab:regimes}
\end{table}

In this paper we measure the claimed-versus-verified gap in the second
regime, and we measure it causally rather than observationally. We build
a harness of five deterministic layers beneath an LLM that only plans,
make each layer independently switchable, and run a controlled ablation
in which the model, the world, the task set and the commit are all held
fixed and the harness is the only thing that varies. Every arm is scored
by the same independent instrument, which is never ablated. Minecraft is
the arena because it is the cheapest environment with the properties
that matter --- unbounded actions, expensive ground truth, long
horizons, and recovery-from-failure as part of the task rather than the
end of it. It is not the target, and we argue for neither the game nor
the harness as artifacts.

\begin{callout}{The design principle}
\centering\normalsize\upshape
\textbf{The claim and the check must not share an author.}
\end{callout}

The intuition behind the design is a separation of authorship. The model
is the right author of a plan, because a plan is a judgment over an
unbounded space. It is the wrong author of a completion check, because a
check is arithmetic over observable state and the model has an interest
in the answer. Generalised: code owns every fact it can compute, and the
model owns only what it cannot. The five layers are that principle
instantiated five times --- supply the arithmetic, gate on it, act on it
deterministically, refuse a claim that contradicts it, and close a task
the world already shows as met.\footnote{The five layers are not five
checks. Two supply facts, one gates, one acts, one refuses and one
terminates; what unifies them is that each moves a decision out of the
model wherever the decision is computable.
\S\ref{sec:layers} states each precisely.} The analogy is not to a
better model but to double-entry bookkeeping: the audit function exists
because the party who benefits from a number should not be the party who
certifies it.

We then study the harness the way a systems paper studies an
architecture, and the results were not what the design predicted. The
whole-harness effect is large and survives every attempt we made to
break it. The per-layer decomposition --- pre-registered, each layer
scored on the failure category it targets --- found \emph{no single layer
responsible}, and the surprise is instructive rather than deflating:
splitting every arm on the same covariate shows the failure has one
trigger and five independent guards, so removing any one of them leaves
it near zero. Capability degrades roughly additively as layers come off;
honesty does not. Alongside this we identify the trigger from stratified
data and then remove it experimentally in a follow-up campaign, which
turns an association measured inside one arm into an intervention.

\medskip
\noindent\textbf{Contributions.}
\begin{enumerate}[label=\textbf{C\arabic*.},leftmargin=2.6em,itemsep=.25em]
\item We draw and then exploit a distinction between \emph{measurement
  regimes}: the setting where checking is cheap, in which false success
  has been measured at scale, and the setting where checking is the
  engineering problem, in which the field still runs on the honor
  system. The two literatures have not met.
\item We run a controlled same-model ablation of a five-layer
  deterministic harness, whole and layer-by-layer, with pre-registered
  endpoints and a minimum effect of interest fixed in advance --- the
  ablation the closest prior work names among its own limitations
  \citep{leni}.
\item We locate the failure's trigger and then remove it
  experimentally, driving false success from 18/24 to 0/24 under a
  decision rule declared before collection; and we report the per-layer
  null as it came out, with the redundancy that explains it.
\item We apply the paper's own standard to the paper's own instrument,
  with a stratified blind human calibration sealed under a published
  hash before any evidence card was rendered, and we release the
  pre-registration, deviations log, ledger, evidence files and analysis
  scripts.
\end{enumerate}

% =====================================================================
\section{Related Work}
\label{sec:related}

\runin{False success where ground truth is cheap.}
A body of recent work measures the gap between claimed and actual
completion directly. The largest characterisation covers 9,876
tau2-bench and 1,879 AppWorld trajectories and reports false success at
45--48\% of failures in single-control domains, rising to 75.8\% among
self-assessing coding trajectories \citep{falsesuccess}. VIGIL
separates the two axes explicitly, distinguishing world completion from
self-termination over 20 models and 1,000 frozen episodes and finding up
to 19.7 points of benchmark difference between systems with
near-identical execution \citep{vigil}. StressWeb shows the
miscalibration persists in clean web environments and grows under
perturbation \citep{stressweb}; the MAST taxonomy places premature
termination at 6.2\% of multi-agent failures \citep{mast}, and Aegis
taxonomises agent--environment failures more broadly \citep{aegis}.
\emph{Different from ours}, all of these are situated where a verdict is
a lookup --- a row, an API response, an exact cell match. They establish
the phenomenon and its prevalence; they cannot speak to the regime in
which producing the verdict is itself the hard part. We claim neither
the phenomenon nor the two-axis distinction as a contribution.

\runin{Self-verification where ground truth is expensive.}
The open-world agent literature operates in exactly that regime and
largely resolves it by delegation to the agent. Voyager's tasks are
self-verified by a critic agent whose approval admits code to a
permanent skill library, so an error compounds on reuse
\citep{voyager}. Luban is built on autonomous embodied verification
\citep{luban}; MineEvolve distils skills from executions its own loop
judged successful \citep{mineevolve}. There is movement toward external
scoring --- MineExplorer uses rule-based milestone evaluators
\citep{mineexplorer} --- but it scores exploration coverage rather than
the claimed-versus-verified gap. \emph{In contrast}, we hold the model
fixed and vary only the machinery that produces the verdict, which turns
the honor system from a background assumption into a manipulable
condition, and we record the claim and the world independently on every
attempt so the gap between them is an observable rather than an
inference.

\runin{Decomposing a scaffold's contribution.}
The closest prior work to our design is a cross-benchmark decomposition
of a production enterprise agent, which separates scaffolding, routing,
specialist models and a verification loop and reports the verification
step's isolated contribution as $+1.5$ points of an $+11.0$ point uplift
\citep{leni}. Two things follow, and both shaped this paper. First,
their verification loop is measured on \emph{task success}, where it
looks marginal --- the same shape as our own secondary endpoint, and a
structural corroboration of our choice to score each layer on the
category it targets instead: overall success dilutes a verification
layer across every run in which its failure mode never arises, and our
own power calculation puts the difference at $n \approx 372$ against
$n \approx 52$ per arm. Second, \emph{different from ours}, that work's
stated limitations include the absence of a controlled same-model
scaffold ablation, alongside vendor evaluation, single scored runs, and
unreleased weights and prompts. We supply the controlled same-model
ablation, repeated replicates, an independent scorer, and open receipts.

\runin{Vocabulary.}
We adopt the field's terms rather than coining our own: \emph{false
success} for what our code calls \texttt{false\_completion}, and on
first use we map our \emph{reached-not-recognized} to VIGIL's
post-attainment drift and our \emph{capability failure} to its missed
execution. A paper that invents a synonym for false success is a paper
nobody searching for false success will find.

% =====================================================================
\section{Method}
\label{sec:method}

\runin{Problem formulation.}
A task $t$ is issued to an agent situated in a world at state $w_0$. The
agent acts and eventually terminates, emitting a terminal report
$c \in \{\text{claimed}, \text{not claimed}\}$ and leaving the world at
$w_1$. Independently, a scorer computes
$s = V(t, w_0, w_1) \in \{\text{success}, \text{fail}\}$ from world
state alone, never from the agent's report. Every criterion in this work
is a \emph{net gain} measured against the task's own start snapshot, so
$V$ is a function of the pair and not of $w_1$ in isolation. The four
cells of $c \times s$ are the unit of analysis, and the cell
$c = \text{claimed} \wedge s = \text{fail}$ is \emph{false success}, the
primary endpoint. Its complement in the claimed row is true completion;
the unclaimed row splits into reached-not-recognized
($s = \text{success}$) and, by cause, capability failure or budget
exhaustion. The estimand is the difference in the rate of a given cell
between two harness configurations, with model, world and task set
fixed.

\runin{Overall architecture.}
The system under test is an LLM that only plans, sitting above five
deterministic layers, each independently switchable by a flag and each
ablated alone in its own arm (Figure~\ref{fig:arch}). One property of
the design matters for interpretation and is easy to misread: the fully
ablated arm does not lose \emph{data}. It still receives the same
\texttt{SELF} / \texttt{INVENTORY} / \texttt{NEARBY} blocks every turn,
at parity with the upstream system. What perception adds is the
\emph{subtraction} --- the delta against the task's start --- and what
the completion check adds is the refusal to accept a claim when that
subtraction disagrees. The intervention is not \emph{give the model the
facts}. It is \emph{do the arithmetic in code, and refuse the claim when
the arithmetic disagrees}.

\begin{figure}[t]
\centering\small
\setlength{\extrarowheight}{3pt}
\begin{tabular}{@{}p{0.20\linewidth}!{\vrule width 1.2pt}p{0.68\linewidth}@{}}
\multicolumn{1}{@{}l!{\vrule width 1.2pt}}{\footnotesize\textsc{model}}
 & \multicolumn{1}{l@{}}{\footnotesize\textsc{deterministic code}}\\
\multicolumn{1}{@{}l!{\vrule width 1.2pt}}{\scriptsize authors the plan and the claim}
 & \multicolumn{1}{l@{}}{\scriptsize authors every check}\\
\noalign{\smallskip\hrule\smallskip}
\scriptsize reads world context &
  \textbf{L1 Perception.} Supplies \texttt{TASK-PROGRESS} --- the delta
  from the task's own start snapshot --- and \texttt{CAPABILITIES}.\\
\scriptsize proposes a step &
  \textbf{L2 Preconditions.} Gates the action on checkable state: is
  this block mineable, is the tool held, is the recipe satisfiable.\\
\scriptsize --- not consulted --- &
  \textbf{L3 Reflexes.} Named deterministic responses to recurring
  situations, run in code rather than proposed by the model.\\
\scriptsize reports ``done'' &
  \textbf{L4 Grounded completion check.} Refuses a finish claim the
  world state does not support.\\
\scriptsize --- not consulted --- &
  \textbf{L5 Deterministic termination.} Closes a task the world shows
  as met, whether or not the agent noticed.\\
\noalign{\smallskip\hrule\smallskip}
\multicolumn{2}{@{}p{0.92\linewidth}@{}}{\scriptsize \textsc{scorer} ---
  never ablated. Reads world state independently over RCON and computes
  the verdict from the start/end pair. The measuring instrument, not a
  component under test.}\\
\end{tabular}
\caption{\textbf{Authorship, not capability.} The five layers sit on the
boundary between what the model decides and what code decides; each is
switchable alone, and the arms of \S\ref{sec:design} are exactly the
settings of these five flags. Two layers supply facts, one gates, one
acts, one refuses and one terminates --- what unifies them is that each
moves a computable decision out of the model. The scorer sits outside
the boundary entirely, which is why it is the one component the ablation
never touches.}
\label{fig:arch}
\end{figure}

\subsection{The five layers}
\label{sec:layers}

Each layer is stated as its input and the decision it removes from the
model. \textbf{Perception} computes the per-task inventory and world
delta and injects it as \texttt{TASK-PROGRESS}, together with a
\texttt{CAPABILITIES} block enumerating what the agent can presently do.
\textbf{Preconditions} refuse an action whose requirements are checkably
absent, converting a doomed attempt into an immediate legible refusal.
\textbf{Reflexes} replace model-proposed responses to recurring
situations with named deterministic ones. \textbf{The grounded
completion check} intercepts a terminal claim and refuses it when the
recomputed delta does not satisfy the criterion. \textbf{Deterministic
termination} is its mirror image: it closes a task whose criterion the
world already satisfies, whether or not the agent recognised it. The
last two are why the taxonomy has two distinct unclaimed-success cells,
and why each layer is scored against a different category in
\S\ref{sec:layerabl}.

\subsection{The scorer, and why it is never ablated}
\label{sec:scorer}

The scorer computes each verdict by one of two paths: an \emph{inventory
delta} against the start snapshot, carrying 328 of the \NPrimary{}
primary rows, or a \emph{block scan} of world geometry over RCON for
structure criteria, carrying the remaining 28. It never reads the
agent's report. An early draft of the design listed the scorer as an
ablation arm; that would have removed the measurement rather than
varying a condition, and it was corrected before any data existed. The
scorer's own trustworthiness is not assumed --- \S\ref{sec:prereg}
states the protocol that tests it and \S\ref{sec:calib} reports the
result.

\subsection{Ablation design}
\label{sec:design}

Seven arms $\times$ 13 tasks $\times$ 4 replicates, within-subject on
tasks. One model (\texttt{claude-sonnet-4-6}), one dedicated world, one
commit --- verified: all \NPrimary{} primary rows carry the same hash.
The arms are the full harness, the fully ablated harness, and the five
single-layer ablations. Arm order and task order are independently
shuffled per replicate, and arm position is recorded per attempt as a
covariate. Three signals are recorded and never collapsed: the scorer's
verdict; the agent's terminal report, logged independently in the task
queue; and the taxonomy derived from their cross-classification. Attempt
disposition, the discard-versus-supersession distinction, and one
disclosed deviation affecting position control are in
Appendix~\ref{app:protocol}.

\subsection{Pre-registration, endpoints, and instrument certification}
\label{sec:prereg}

\runin{Endpoints and decision rules.}
Endpoints were registered before collection. Each layer is scored on the
taxonomy category it targets rather than on overall success, with a
minimum effect of interest of 15 points declared in advance and applied
in code as a verdict, rather than left in prose for a reader to apply by
eye \citep{kohavi2020}. Inference is by cluster bootstrap with the task
as the resampling unit and \Resamples{} resamples, Wilson intervals for
single proportions, McNemar exact for the within-attempt claim/verify
contrast, and Holm--Bonferroni across the five-layer family.

\runin{Certifying the instrument.}
A paper whose subject is that self-report cannot be trusted owes the
same treatment to its own scorer. The question is not whether a blind
human agrees with the scorer on average, but whether one agrees \emph{on
the verdict the paper depends on} --- which is rare by construction. The
protocol therefore samples by verdict rather than at random: thirty
attempts drawn from this campaign's own primary set, stratified by
scorer verdict and by which scoring path produced it, oversampling the
false-success cell roughly fifteen-fold.

\begin{callout}{Registered before any card was rendered}
The sampling rule, allocation, RNG seed and decision rule were committed
first. The stratum-annotated manifest is the answer key, so it was
sealed behind a SHA-256 and withheld, leaving only the hash and a
stratum-blind presentation order; the draw re-derives from the seed and
verifies against the seal.

\smallskip
The rule fixed in advance: \textbf{report per cell and never pool},
because the sample is deliberately non-representative. Below 90\%
agreement in a cell, apply the measured error rate as a correction to
every rate depending on it and report a corrected instrument --- not a
limitation of the labels, and not grounds to fall back on an earlier
figure. Abstentions are recorded as abstentions and counted as neither
agreement nor disagreement.
\end{callout}

Because stratified sampling by verdict fixes the column margins of the
$2 \times 2$, the per-cell rates are predictive values rather than
sensitivity and specificity. They answer the question a reader of a rate
actually has --- \emph{given that the scorer said this, how often is it
right} --- and cannot be read as a single accuracy figure. Seed,
allocation, seal hash, verification command and prevalence reweighting
are in Appendix~\ref{app:calib}.

% =====================================================================
\section{Experiments}
\label{sec:exp}

\runin{Environment and tasks.}
All attempts run against a dedicated Minecraft server instance reserved
for evaluation and isolated from any production world, with the scorer
reading state over RCON. The task set is 13 tasks spanning gathering,
crafting, tool progression and one structure criterion, each expressed
as a net gain from the attempt's own start snapshot. Every attempt
begins from a fixed starting kit built from the observed campaign median
of each consumable, and the kit is deliberately unable to satisfy any
criterion on its own: ``$+16$ cobblestone this run'' still requires
mining 16 with 96 already in the bag. That intent is stated in the
frozen runner, written before collection began, and it is the design
decision \S\ref{sec:mech} turns out to depend on. Task definitions, kit
composition and per-arm trial counts are in Appendix~\ref{app:eval}.

\runin{Implementation details.}
One model, \texttt{claude-sonnet-4-6}, at fixed decoding settings across
all arms; one commit across all \NPrimary{} primary rows, verified by
hash. Each attempt runs under a fixed step and wall-clock budget, with
budget exhaustion recorded as its own taxonomy category rather than
folded into failure. Every attempt writes one append-only ledger row
plus an evidence file; every taxonomy label, derived metric and
statistic is recomputed downstream from those raw receipts, and every
table in this manuscript is generated from the ledger rather than
transcribed.

\subsection{Instrument certification}
\label{sec:calib}

\runin{Why the earlier figure does not certify this campaign.}
An earlier calibration reported 39 of 39 blind labels in agreement,
Wilson $[91.0, 100.0]$. That figure is accurate about the population it
was drawn from and it does not certify this campaign. Cards had been
issued in glob order with no stratification, and regenerating the
composition shows why that matters: 34 of the 39 labels are true
completions, 3 are capability failures, and 2 are false successes. None
of the 39 judges an attempt from this campaign at all. The scorer's
false-success verdict --- what every headline number here rests on ---
was human-validated on two attempts, Wilson $[34.2, 100.0]$. Aggregate
agreement is the wrong statistic for a rare-cell endpoint, and that is
what motivated the stratified protocol of \S\ref{sec:prereg}.

\begin{table}[t]
\centering\small
\begin{tabular}{@{}lrrrr}
\toprule
cell & in primary set & labelled & agreement & Wilson 95\%\\
\midrule
\inputrows{tables/calibration.tex}
\bottomrule
\end{tabular}
\caption{Blind human agreement, by scorer verdict cell. Reported per
cell and never pooled, as pre-registered.}
\label{tab:calib}
\end{table}

\runin{Result.}
The pre-declared rule returns a pass on all three cells. The endpoint
cell moves from $n=2$ to $n=21$ at 20/21 $=$ 95.2\%, and on the
inventory-delta path that scores 328 of the \NPrimary{} primary rows it
is 17/17 with a lower bound of 81.6\%. No correction to any downstream
rate is triggered.

\runin{The one disagreement, and what it costs us.}
On a structure task the agent claimed completion, the block scan found
16 of 25 blocks and scored a false success, and blind human judgement
called it a genuine success. It is counted as a disagreement in every
figure above; the pre-registration forbids revising it now that the
verdict is visible. The diagnostic changes no number and is worth
recording: the evidence card showed a $-25$ cobblestone delta and
nothing about geometry, so the labeller could see that twenty-five
blocks were placed but not that they failed to form a solid
$5 \times 5$. \textbf{On structure criteria the human is the weaker
instrument, not the stronger one} --- which inverts the assumption blind
calibration is built on, that the human sees everything and can catch a
criterion the rule mis-encodes. That assumption holds for an inventory
delta and fails for a shape. We therefore do not claim the block-scan
path is human-calibrated; its warrant is engineering validation against
ground truth placed over RCON, which for a deterministic geometric check
is meaningful evidence and is not what a calibration claims.

\runin{A counter-exhibit.}
Where the label came from the agent's own claim rather than from the
scorer, blind human judgement disagreed every time --- 0 of 5, across
both calibration eras. That is the phenomenon under study, measured on
the instrument that is supposed to be the alternative to it.

\subsection{Main result: whole-harness ablation}
\label{sec:main}

\begin{callout}{Primary and co-primary endpoints}
\upshape
The two arms claim success at statistically indistinguishable rates and
deliver less than half as much. The harness does not make the agent
claim less --- it makes the claims true.

\smallskip
\noindent\begin{tabular}{@{}ll@{}}
claimed        & full harness \textbf{\OnClaim{}} $\cdot$ no harness \textbf{\OffClaim{}}\\
verified       & full harness \textbf{\OnSucc{}} $\cdot$ no harness \textbf{\OffSucc{}}\\
task success   & \textbf{\SuccDiff{}} \quad CI \SuccCI{} \quad $p \le 0.0002$\\
false success  & \textbf{\FCDiff{}} \quad CI \FCCI{} \quad $p \le 0.0002$\\
\end{tabular}
\end{callout}

\begin{table}[t]
\centering\small
\setlength{\tabcolsep}{4pt}
\begin{tabular}{@{}lrrcrrrrr@{}}
\toprule
arm & $n$ & verified & 95\% CI & claimed & $b$ & $c$ & gap & false succ.\\
\midrule
\inputrows{tables/arms.tex}
\bottomrule
\end{tabular}
\caption{Every arm. $b$ = claimed but not verified; $c$ = verified but
never claimed. All rates in percent; gap $=(b-c)/n$.}
\label{tab:arms}
\end{table}

\runin{Direction, not just magnitude.}
McNemar's exact test on claimed-versus-verified within the same attempts
gives $p = 2.4 \times 10^{-7}$ for the ablated arm, and every one of its
23 discordant pairs runs the same way --- the agent claiming success it
did not have, never the reverse. Under the full harness there are two
discordant pairs and the gap is not statistically detectable. A
symmetric error would be noise; a one-directional one is the phenomenon.

\runin{The say-do gap is not the endpoint.}
Two arms show why they cannot be used interchangeably. The
\emph{$-$~reflexes} arm posts a gap of exactly zero while holding one
false success and one unrecognized success --- two errors cancelling,
not an honest arm. The gap is an aggregate; false success is the
endpoint, and Table~\ref{tab:arms} reports $b$ and $c$ separately for
exactly this reason.

\runin{Arms fail in different kinds.}
A success rate alone hides this. The fully ablated arm fails mostly by
\emph{claiming} --- 23 false successes against 9 capability failures.
Every single-layer arm fails mostly by \emph{not managing}, with
capability failure dominant and false success in low single digits. The
full taxonomy by arm is in Appendix~\ref{app:eval}.

\subsection{Per-layer ablation}
\label{sec:layerabl}

The per-layer decomposition was the campaign's fourth hypothesis and
\textbf{it is not supported}. Scored on the category each layer targets,
with Holm correction across the five comparisons, no layer moved its
target failure mode by the pre-declared 15 points.

\begin{table}[t]
\centering\footnotesize
\setlength{\tabcolsep}{3pt}
% Seven columns do not fit the text block at any sane font size; scaled to
% width rather than shrunk further, so the whole table stays one size.
\resizebox{\textwidth}{!}{%
\begin{tabular}{@{}llrcrlc}
\toprule
layer removed & target category & rise & 95\% CI & Holm & verdict &
  platform-drop\\
\midrule
\inputrows{tables/layers.tex}
\bottomrule
\end{tabular}}
\caption{Per-layer ablation against the full harness, each scored on its
own endpoint. The final column repeats the estimate with the non-delta
platform attempt removed.}
\label{tab:layers}
\end{table}

\runin{The near miss, and why we do not report it as the mechanism.}
The completion check is the closest and points the predicted way, but it
neither survives multiplicity nor clears the minimum effect of interest
--- and its unadjusted interval is carried by a single task. Dropping
the one criterion decided by the block-scan path, it falls to
$+9.8$pp, $[-1.5, +21.4]$, spanning zero. The one per-layer result that
looked nearly detectable is also the one leaning hardest on the
least-certified instrument in the rig.

\runin{Attribution.}
Splitting every arm on the covariate identified in \S\ref{sec:mech}
explains the null directly (Table~\ref{tab:strata}). Removing any single
layer leaves false success at or near zero; only removing all five
produces 75\%. We attribute the null to \emph{redundancy} --- the
failure has one trigger and five independent guards, any one of which is
sufficient, so no single-layer ablation could have reached a 15-point
threshold. \S\ref{sec:disc} gives the arithmetic and states what the
reading does and does not establish.

\begin{table}[t]
\centering\small
\begin{tabular}{@{}lrr@{}}
\toprule
arm & kit supplies the target & kit empty of it\\
\midrule
full harness                & 1/24 = 4\%   & 0/24 = 0\%\\
no harness                  & 18/24 = 75\% & 4/23 = 17\%\\
$-$ perception              & 0/24 = 0\%   & 0/24 = 0\%\\
$-$ preconditions           & 0/24 = 0\%   & 0/24 = 0\%\\
$-$ reflexes                & 0/23 = 0\%   & 1/24 = 4\%\\
$-$ completion check        & 2/21 = 10\%  & 3/21 = 14\%\\
$-$ determ. termination     & 0/24 = 0\%   & 0/24 = 0\%\\
\bottomrule
\end{tabular}
\caption{False success by arm and starting-stock stratum. Regenerate
with \texttt{python3 eval/stock\_strata.py}.}
\label{tab:strata}
\end{table}

\subsection{What the false successes are, and a direct manipulation}
\label{sec:mech}

\runin{The trigger.}
The unharnessed agent treated the starting bag as the answer. Among its
false successes, 18 of 22 began holding at least the required quantity
of the target item, against 5 of 14 among its true completions --- odds
ratio 8.1, Fisher exact $p = 0.0112$. It is not partial credit: 21 of
the 22 produced literally zero gain on the target item, at a median of 2
steps and 10 seconds. The agent read an absolute quantity as if it were
the required gain. The denominator is 22 rather than the arm's 23
because the stratum rule --- \emph{did this attempt begin holding at
least $N$ of its own target item} --- requires a countable target, so
the structure-criterion attempts are absent from every stratified row
here, and one of the arm's 23 false successes is among them.

\begin{callout}{Representative trajectories, ablated arm}
\upshape\ttfamily\footnotesize
needed +12 stick \quad had 96 \quad gained 0 \quad $\rightarrow$ claimed done\\
needed +64 cobblestone \quad had 96 \quad gained 0 \quad $\rightarrow$ claimed done\\
needed +16 oak\_planks \quad had 48 \quad gained 0 \quad $\rightarrow$ claimed done\\
needed +1 stone\_pickaxe \quad had 1 \quad gained 0 \quad $\rightarrow$ claimed done
\end{callout}

\runin{The obvious objection, and the matched test.}
If the agent already holds 96 sticks and is asked for 12 more, refusing
to grind is arguably sensible, and scoring it a failure would penalise
good judgement rather than dishonesty. The test is matched --- same
tasks, same starting stock, only the harness differs. The same model, on
the same task, holding the same 96 sticks, crafted 12 more four times
out of four when the harness was on. Across the six stock-satisfied
tasks: 20/24 achieved the required net gain with the harness, 5/24
without. The objection does identify the correct boundary, and the
taxonomy respects it: had the agent answered ``I already hold 96 sticks
--- do you still want 12 more?'', that is a question or a cancel, scored
as neither a true nor a false completion. What is scored is asserting
completion of an action that did not occur.

\runin{The manipulation.}
Everything above is an association measured within one arm; the matched
test is causal for the harness and not for the stock. A follow-up
campaign supplies that manipulation. The six stock-satisfied tasks were
re-run fully ablated with a kit that drops \emph{only} that task's
target item and keeps every enabler, so the work required is unchanged
and the sole variable is whether the goal was already satisfied at the
start.

\begin{callout}{Pre-declared: $\le$33\% confirms $\cdot$ 34--49\% inconclusive $\cdot$ $\ge$50\% refutes}
\upshape
Removing the target item from the kit drove false success from
\textbf{18/24 $=$ 75\%} to \textbf{0/24 $=$ 0\%}.

\smallskip
\noindent\ttfamily\footnotesize
Fisher exact, two-sided p = 3.7e-08 $\cdot$ rule returns CONFIRMED at 0\%
\end{callout}

Twenty of the 24 succeeded, most overshooting the target. The four that
failed did so openly --- three timeouts and a cancel --- and \emph{not
one of them claimed completion}. That is the distinction the endpoint
measures: with an empty bag the agent either did the work or failed
without reporting success.

\runin{Why the level does not travel, and the contrast does.}
The \OffFC{} headline is the average of two very different cells, and
only one of them generalises.

\begin{table}[t]
\centering\small
\begin{tabular}{@{}lrrr@{}}
\toprule
stratum & no harness & full harness & Fisher exact\\
\midrule
kit supplies the target item & 18/24 = 75\% & 1/24 = 4\% & $5.7 \times 10^{-7}$\\
kit empty of it              & 4/23 = 17\%  & 0/24 = 0\% & $0.0496$\\
\bottomrule
\end{tabular}
\caption{False success by starting stock. Both arms received the
identical kit, so each stratum is a fair comparison.}
\label{tab:stock}
\end{table}

The harness effect survives in both strata, so \emph{the contrast} is
the transferable finding. The level is a property of our kit: 6 of 13
tasks fall in the first stratum by design, because the kit was built
from the observed campaign median of each consumable. In an
empty-inventory setting the ablated arm would false-complete far less
often. We report this beside the headline rather than as a limitation
three sections later, because a qualification a reader meets that late
is one most readers never meet. Two honesties: the empty-bag contrast is
the weaker of the two, four events against zero at $p = 0.0496$; and the
residual 17\% is a second, smaller failure mode that starting stock does
not explain.

\subsection{Robustness}
\label{sec:robust}

We attempted to break the headline result from five directions and could
not. Full numbers are in Appendix~\ref{app:robust}.

\begin{itemize}[leftmargin=1.3em,itemsep=.3em]
\item \textbf{No replicate drives it.} Leave-one-replicate-out gives
  success differences from $+29.8$ to $+43.6$pp, every interval
  excluding zero.
\item \textbf{The exclusion machinery is not load-bearing.} The
  maximally adversarial analysis --- discarding nothing, including all
  63 faults and all 77 superseded rows --- still gives
  $+31.3$pp $[+15.1, +48.1]$ on success and $-37.0$pp
  $[-55.4, -18.9]$ on false success.
\item \textbf{Both sensitivity drops enlarge the effect.} Dropping the
  block-scan task, or dropping the 8 zero-step timeouts, moves both
  endpoints further from zero. They are reported as collected; removing
  rows after seeing which way it moves the result is what
  pre-registration exists to forbid, and keeping them costs us effect
  size rather than manufacturing it.
\item \textbf{The clustering choice is the conservative one.}
  Clustering on replicate rather than task gives narrower intervals.
\item \textbf{Position does not explain it.} With the position reset
  broken (Appendix~\ref{app:protocol}), we restrict to the two
  replicates in which every arm ran in the same region of the world:
  the stock-satisfied contrast is 12/12 versus 0/12 there. This is what
  would be expected of a failure mode that consists of not moving ---
  21 of 22 false successes produced zero net gain at a median of 2
  steps. Terrain does not enter into it.
\end{itemize}

\runin{A design-based check.}
Because the cluster bootstrap resamples only 13 task clusters, against
conventional guidance of 30--50 or more, we add a test that assumes
nothing about independence across tasks. Every arm runs the same task
set within a replicate, so each (replicate $\times$ task) yields a
matched pair, and under the null the two labels are exchangeable --- the
randomization analogue of the McNemar above.

\begin{table}[t]
\centering\small
\begin{tabular}{@{}lrrr@{}}
\toprule
endpoint & observed & permutation $p$ & bootstrap $p$\\
\midrule
task success  & $+39.2$pp & $0.0003$        & $\le 0.0002$\\
false success & $-41.2$pp & $\le 0.0001$    & $\le 0.0002$\\
\bottomrule
\end{tabular}
\caption{Permutation test over 51 matched pairs, \Resamples{}
sign-flips. The two methods are wrong in different ways, so agreement is
evidence rather than restatement.}
\label{tab:perm}
\end{table}

\subsection{Cost}
\label{sec:cost}

Per run the two whole-harness arms are indistinguishable --- 89,632
against 90,375 input tokens, a difference of 0.8\%. Per unit of work
that actually happened, the harness costs 133,167 tokens per verified
success against 307,274, a factor of $2.3\times$. We flag what that is
and is not. Because per-run cost is indistinguishable, the ratio is
arithmetically the success ratio restated, not independent evidence. It
is worth reporting for one reason: tokens-per-run and
tokens-per-verified-success \emph{rank the arms differently}, and only
one of them measures something a user would pay for. That is this
paper's own thesis applied to its own cost table. Per-arm cost for all
seven arms is in Appendix~\ref{app:eval}.

% =====================================================================
\section{Discussion}
\label{sec:disc}

\runin{Capability degrades additively; honesty does not.}
The redundancy reading of \S\ref{sec:layerabl} also appears in the
arithmetic. On false success the five single-layer effects sum to
$+11.7$pp against a whole-harness rise of $+41.3$pp --- an interaction
gap of $+29.6$pp. On task success the same decomposition is nearly
additive, $+34.9$ against $+37.9$. Two endpoints, one design, opposite
additivity: removing guards degrades what the agent can do a little at a
time, and degrades whether its report can be trusted only when the last
guard comes off. If that holds generally it has a practical consequence
--- a scaffold's honesty contribution cannot be estimated by ablating
its parts one at a time, which is the standard way scaffolds are
evaluated.

\begin{callout}{What this is not}
Single-layer ablations cannot identify an interaction. The redundancy
reading is a hypothesis the data suggests and does not demonstrate; a
$2^5$ factorial, or at minimum a joint perception-plus-completion-check
ablation, is the test that would settle it. We report the null as it
came out and the reading as exploratory.
\end{callout}

\runin{Limitations.}
Stated as consequences, not caveats.

\begin{itemize}[leftmargin=1.3em,itemsep=.3em]
\item \textbf{No layer-localisation claim is available.} The campaign
  was sized for a 15-point per-layer effect and none was found. Whether
  the layers genuinely share the work or the design is underpowered for
  the true effect sizes is not resolved here.
\item \textbf{The scorer is certified on one path and not the other.}
  The inventory-delta path carries 328 of \NPrimary{} primary rows and
  is certified at 17/17. The 28 block-scan rows rest on engineering
  validation; a labeller cannot check geometry from an inventory delta,
  so this gap is a property of the evidence card and is not closed by
  drawing more cards of the same kind.
\item \textbf{The absolute false-success rate is kit-conditional}
  (\S\ref{sec:mech}). The contrast transfers; the level does not.
\item \textbf{A 15.0\% infrastructure fault rate is high} for a
  campaign this size, and every fault is a run the rig lost rather than
  a result.
\item \textbf{Unequal $n$.} Three arms hold fewer than \OnN{} attempts,
  and the arm carrying the largest per-layer effect holds 46 --- the
  smallest, most truncated, and largest-effect arm are the same arm.
\item \textbf{Position control rests on randomization alone,} because
  the reset teleport never fired (Appendix~\ref{app:protocol}).
\item \textbf{Family-wise error is controlled within the per-layer
  family only.} The three primary hypotheses are three tests on two
  arms of largely overlapping data. The effects are large enough that
  this changes nothing, and it should be stated rather than left
  implicit.
\item \textbf{One model, one world, one task family.} Nothing here
  separates the harness effect from properties of this model or this
  environment. The ablated arm is this fork with every layer switched
  off --- the correct control for attribution, and not a claim about any
  upstream system.
\end{itemize}

% =====================================================================
\section{Conclusion}
\label{sec:concl}

We have measured the gap between what an open-world LLM agent claims and
what it delivers, in the regime where establishing ground truth is
itself the engineering problem rather than a lookup. Holding the model,
the world, the task set and the commit fixed and varying only a
five-layer deterministic harness, we find two arms that claim success at
indistinguishable rates and differ by \SuccDiff{} in verified success,
with ablation raising false success from \OnFC{} to \OffFC{} and every
discordant attempt running the same direction. The mechanism is
identifiable and removable: false success concentrates where the
starting kit already satisfied the goal, and stripping only that item
eliminates it entirely under a rule fixed in advance. The pre-registered
per-layer decomposition found no single layer responsible, and the
stratified data explains it as redundancy --- capability degrades
additively as guards are removed and trustworthiness does not, which
suggests a scaffold's honesty contribution cannot be estimated one
component at a time. Applying the same standard to our own instrument, a
sealed stratified blind pass certifies the verdict the paper rests on
and declines to certify the path it cannot. What transfers is not a rate
but a contrast, and not a harness but a principle: the claim and the
check must not share an author.

% =====================================================================
\bibliographystyle{plainnat}
\bibliography{refs}

% =====================================================================
\clearpage
\appendix
\section{Protocol, disposition and deviations}
\label{app:protocol}

\runin{Disposition of every attempt logged.}
No attempt is ever deleted; discards are recorded against immutable rows
rather than by editing them.

\begin{table}[h]
\centering\small
\begin{tabular}{@{}lr@{}}
\toprule
 & attempts\\
\midrule
logged                                                    & 496\\
primary set                                               & \NPrimary{}\\
discarded --- infrastructure faults                       & 63 (15.0\%)\\
superseded by a declared whole-replicate re-run           & 77\\
never logged (bot death, retry exhaustion, collisions)    & 8\\
\bottomrule
\end{tabular}
\caption{Disposition of every attempt logged.}
\label{tab:disposition}
\end{table}

\runin{Faults and supersessions are counted separately and never summed.}
A discard is an infrastructure fault, meaning the attempt was never
evidence; a supersession is valid data replaced wholesale by a declared
re-run. Reporting them together produced a 28.2\% ``fault rate'' against
a true 15.0\%, inventing a reliability problem we did not have. Both
leave the primary set.

\begin{callout}{Disclosed deviation}
The per-attempt state reset restored inventory as specified --- 484 of
496 attempts (97.6\%) begin on an inventory exactly equal to the fixed
kit --- but its teleport was \emph{silently rejected}, so every attempt
resumed where the previous one ended.

\smallskip
Root cause is a defect class rather than a typo: the RCON client returns
each command's response and never inspects it, and Minecraft answers a
refused command with an error string over a normal response, so a
declined teleport is indistinguishable from one that worked.
Consequently arm-position control rests entirely on the per-replicate
randomization of arm order, not on the reset; \S\ref{sec:robust} reports
the position-matched analysis this forces. Repair of the response
checking is an open item, scheduled before any future campaign rather
than during one.
\end{callout}

\runin{Pre-registration and the deviations log.}
The frozen protocol, the hypothesis set, the endpoint assignments, the
minimum effect of interest, and an append-only deviations log are
released in full. Every deviation above appears in that log with the
date it was discovered and the analysis it changed.

% ---------------------------------------------------------------------
\section{Evaluation details}
\label{app:eval}

\runin{Task set.}
13 tasks spanning gathering, crafting, tool progression and one
structure criterion, each expressed as a required net gain from the
attempt's own start snapshot. One further platform-construction task is
excluded from the primary set as a non-delta criterion; it is the
excluded row referenced in the ledger, and no give-tasks exist in this
task set.

\runin{Starting kit.}
Fixed across every arm and every attempt, and identical between the
compared arms in every stratified table. Each consumable quantity is the
observed median for that item across an earlier campaign, which is what
places 6 of the 13 tasks in the stock-satisfied stratum of
\S\ref{sec:mech}. The kit cannot satisfy any criterion on its own,
because every criterion is a net gain.

\runin{Trial counts.}
Nominally 4 replicates $\times$ 13 tasks $=$ \OnN{} attempts per arm.
Realised counts are in Table~\ref{tab:arms}: four arms at 52, two at 51,
and the completion-check arm at 46.

\begin{table}[h]
\centering\small
\setlength{\tabcolsep}{4pt}
\begin{tabular}{@{}lrrrrrr@{}}
\toprule
arm & $n$ & true compl. & false succ. & reached, not recog. &
  capability & budget\\
\midrule
\inputrows{tables/taxonomy.tex}
\bottomrule
\end{tabular}
\caption{Outcome taxonomy by arm.}
\label{tab:taxonomy}
\end{table}

\begin{table}[h]
\centering\small
\begin{tabular}{@{}lrrr@{}}
\toprule
arm & input tokens / run & seconds / run & tokens / verified success\\
\midrule
\inputrows{tables/cost.tex}
\bottomrule
\end{tabular}
\caption{Cost per arm.}
\label{tab:cost}
\end{table}

\runin{Platform caveat.}
Per-layer intervals are reported with and without the non-delta platform
attempt, since a single task can carry a per-layer interval at these
sample sizes; the completion-check arm's drop from $+13.5$ to $+9.8$pp
under that removal is the case in point (\S\ref{sec:layerabl}).

% ---------------------------------------------------------------------
\section{Calibration details}
\label{app:calib}

\runin{Draw and seal.}
Thirty attempts were drawn from the primary set with a fixed RNG seed,
stratified by scorer verdict cell and by scoring path, with the
false-success cell oversampled roughly fifteen-fold. The
stratum-annotated manifest is the answer key and was sealed under a
published SHA-256 before any evidence card was rendered; labellers
received a stratum-blind presentation order only. The draw re-derives
from the seed and verifies against the seal, and the verification
command is released alongside the manifest. One attempt whose identifier
mapped to two ledger rows was dropped from the draw rather than resolved
after the fact.

\runin{Abstentions and coverage.}
One block-scan attempt is recorded as an abstention and counted as
neither agreement nor disagreement, so the denominator remains
auditable. One stratum, 8 attempts, carries no label at all; strata
coverage is 95.2\% of the primary set.

\runin{Prevalence reweighting.}
Reweighting the per-cell rates by the primary set's composition gives
99.2\% overall with a conservative bound of 52.9\%. The bound is
computed from weighted Wilson lower bounds renormalised over covered
strata, not from a normal approximation: three of these strata came back
at 100\%, where the Wald variance is exactly zero and a
normal-approximation interval would collapse to a spuriously tight
point. The bound is weak because the label budget was deliberately spent
on the cell that carries the paper rather than spread evenly, and it is
reported as weak. \textbf{The strong claim is per-cell and specifically
about false success.}

\runin{Why predictive values and not sensitivity.}
Sampling by the scorer's verdict fixes the column margins of the
$2 \times 2$, which yields positive and negative predictive values
conditional on the verdict and does not identify sensitivity or
specificity. This is the quantity a reader of a rate needs, and it is
the reason the pre-registration forbids pooling the cells into a single
accuracy figure.

\runin{Label eras.}
Three eras are tracked separately in the label store --- a pilot era,
the earlier unstratified confirmatory era whose composition is reported
in \S\ref{sec:calib}, and this stratified era --- because assigning an
era by label timestamp would have pooled the stratified oversample into
the earlier 39/39 anchor and inflated it. The anchor is preserved
uncontaminated at 39/39.

% ---------------------------------------------------------------------
\section{Robustness details}
\label{app:robust}

\runin{Leave-one-replicate-out.}
Four refits, each omitting one replicate: success differences $+29.8$ to
$+43.6$pp, every interval excluding zero.

\runin{Maximally adversarial inclusion.}
Discarding nothing --- all 63 infrastructure faults and all 77
superseded rows re-admitted --- gives $+31.3$pp $[+15.1, +48.1]$ on task
success and $-37.0$pp $[-55.4, -18.9]$ on false success. This is the
analysis a sceptic would demand, and it is reported rather than the
favourable one.

\runin{Sensitivity drops.}
Dropping the block-scan-decided task, and separately dropping the 8
zero-step timeouts, each move both endpoints further from zero. Both are
reported as collected.

\runin{Clustering.}
The task is the resampling unit throughout, at 13 clusters. Clustering
on replicate instead gives narrower intervals, so the reported choice is
the conservative one. The permutation test of \S\ref{sec:robust} is the
design-based complement that assumes nothing about cross-task
independence.

\runin{Position.}
Restricting to the two replicates in which every arm ran in the same
region of the world gives a stock-satisfied contrast of 12/12 versus
0/12. A failure mode consisting of not moving cannot be a terrain
artifact.

\runin{P-value floors.}
The cluster bootstrap uses \Resamples{} resamples, and a two-sided
bootstrap $p$ is floored at $2/(B+1)$, so the smallest reportable value
is $p \le 0.0002$ rather than the $1/(B+1)$ a one-sided reading would
suggest. Fisher and McNemar values are exact and reported in scientific
notation below $10^{-4}$.

\runin{Deferred.}
A mixed-effects (GLMM) cross-check alongside the permutation test is an
open item; the analysis stack is deliberately standard-library-only, so
it requires a dependency the campaign did not carry.

% ---------------------------------------------------------------------
\section{Artifacts and disclosure}
\label{app:artifacts}

\runin{Released.}
The pre-registration with its full deviations log, the frozen protocol,
the abort log, the ledger, the evidence files, the analysis scripts, and
the sealed calibration draw with its verification command. Every attempt
is one append-only ledger row plus an evidence file; raw receipts are
immutable, and every taxonomy label, derived metric, statistic and
manuscript table is recomputed downstream from those receipts rather
than transcribed.

\begin{callout}{AI assistance}
This work was carried out with substantial AI assistance, including the
literature audit, analysis code, and drafting. The cited works were
machine-screened against the arXiv API --- all ten identifiers resolve
and every title matches what is attributed to it here --- and three
specific quantitative claims remain pending human verification against
the source PDFs before submission. A machine confirming a machine's
citations rules out fabricated identifiers and wrong titles; it does not
discharge the obligation, and the sign-off is human.
\end{callout}

\runin{Open items before submission.}
Human verification of three attributed claims; a joint
perception-plus-completion-check ablation to test the redundancy
hypothesis directly; blind labels on structure-scored attempts, which
requires an evidence card that can convey geometry; a mixed-effects
cross-check alongside the permutation test; and repair of the RCON
response checking that caused the position-reset failure.

\end{document}
